\scanpage{94}
% archive foliation: 93
et $\mathcal{L}\otimes\underline{\omega}$, induisant $-2i$ et $\operatorname{id}$. \unsure{On a} un iso
\[
\mathcal{L}\underset{\gamma^{-1}\otimes\operatorname{id}_{\omega^{-1}}}{\simeq}P^1_{X/S}(\underline{\Omega})\omega^{-1}\simeq P^1_{X/S}(\omega)\omega^{-1}
\]
\deleted{\unsure{Cor.}}
induisant un diagramme d'extensions
\diagram{\begin{tikzcd}
0\arrow[r] & \underline{\omega}\arrow[r]\arrow[d,"\operatorname{id}"] & \mathcal{L}\arrow[r]\arrow[d] & \mathcal{O}_X\arrow[r]\arrow[d,"{-2\operatorname{id}}"{draw,circle}] & 0\\
0\arrow[r] & \underline{\omega}\arrow[r] & P^1_{X/S}(\omega)\omega^{-1}\arrow[r] & \underline{\mathcal{O}}_X\arrow[r] & 0
\end{tikzcd}}
\margin{\textcircled{$\alpha$}}
ou aussi, le multipliant par $-\frac12$
\diagram{\begin{tikzcd}
0\arrow[r] & \underline{\omega}\arrow[r]\arrow[d,"{-\frac12\operatorname{id}_{\omega}}"{draw,circle}] & \mathcal{L}\arrow[r]\arrow[d] & \mathcal{O}_X\arrow[r]\arrow[d,"{\operatorname{id}_{\mathcal{O}_X}}"] & 0\\
0\arrow[r] & \underline{\omega}\arrow[r] & P^1_{X/S}(\omega)\omega^{-1}\arrow[r] & \mathcal{O}_X\arrow[r] & 0
\end{tikzcd}}
\margin{\textcircled{$\unsure{\alpha'|\beta}$}}

(Donc le torseur sous $\underline{\omega}\simeq\Omega$ \unsure{définissant} l'ext. $\mathcal{L}$ est déduit du torseur \inserted{\unsure{sous $\omega$}} défini \inserted{par l'ext.} $(P^1_{X/S}(\underline{\omega})\underline{\omega}^{-1}$ de $\mathcal{O}_X$ par $\underline{\omega}$ (ou $P^1_{X/S}(\omega)$ de $\underline{\omega}$ par $\underline{\omega}^2)$ \uline{par multiplication par $-2$}.

Utilisant ces hom. comme identifications,
$c|\underset{\substack{\Vert\\\operatorname{Fil}_1\mathcal{H}\\\Vert\\\mathcal{L}^{\unsure{\mathrm{ab}}}}}{\underline{\Omega}\simeq\underline{\omega}}$ s'identifie à l'op. diff. universel,
\[
d^1_{\underline{\omega},X/S}\quad
\underset{\substack{\Vert\\\omega}}{\Omega}\xrightarrow{c}\mathcal{L}\otimes\omega\simeq P^1(\omega)
\]

\uline{NB} On utilisera \unsure{l'iso} comme identification et on se borne aux conn. de $\mathcal{H}$, donc $c:\mathcal{H}\longrightarrow\mathcal{H}\otimes\omega$, dont la restriction à $\underline{\Omega}\simeq\underline{\omega}=\unsure{\mathcal{L}_{\omega}}$ est le $d^1_{\underline{\omega},X/S}$.
\deleted{Ils sont déterminés \unsure{par une}} \deleted{\unsure{famille}} \inserted{\unsure{section}} \unsure{donnée à} $\underline{\Omega}\otimes\omega\simeq\underline{\omega}^{\otimes2}$.
